% ===================================================================
%  TABELLA N-RO 12 — Le station. — Le ferrovias. — Le naves.
%  Traduction 2026 d'apres texto/io, controlee sur texto/fr et
%  texto/es. Consigne : voir texto/ia/10-tabella-01.tex.
% ===================================================================

%%K t12-apar-1 apar
\VUsaut{4.0mm}
\VUtitre{15.0pt}{60}{TABELLA N\textsuperscript{o} 12}
\VUsaut{2.4mm}
\VUpk{11.4pt}{40}{Le station. --- Le ferrovias. --- Le naves.}
\VUsaut{3.4mm}
\textsc{Nota.} --- \textit{Le horarios usate in cata nation es diverse, nos dunque
usa como exemplo le horas del} \VUgras{tabella francese.}

%%K t12-01-1 p
1. --- Io justo viagiava tra Germania. Ante le partita, io comprava un
\VUgras{indicator} \textsuperscript{(15)} pro saper le hora de partita del
trainos. Post constatar que le plus opportun traino va partir de Paris, a novem
horas del vespere, e arrivar in Strasbourg a un hora dece-tres minutas, io
preparava mi viage, e le die sequente, io me faceva conducer al station.

%%K t12-02-1 p
2. --- Quando io entrava le grande hall de partita, detra un vetule
\VUgras{sacerdote} \textsuperscript{(2)} rural, qui portava in le mano su
\VUgras{sacco de viage} \textsuperscript{(3)}, multe \VUgras{viageros}
\textsuperscript{(1)} ja habeva arrivate.

%%K t12-02-2 p
Proque io voleva inviar un \VUgras{dispatcho (telegramma)}
\textsuperscript{(5)}, io me faceva indicar per un \VUgras{controlator}
\textsuperscript{(6)} \VUgras{le officio telegraphic} \textsuperscript{(4)}.

%%K t12-03-1 p
3. --- Ante passar in le \VUgras{sala de attender} \textsuperscript{(7)} io iva
provider me de un \VUgras{billet} \textsuperscript{(9)} de ir e retorno.

%%K t12-04-1 p
4. --- Io non attendeva longemente al \VUgras{guichet} \textsuperscript{(8)},
proque pauc personas esseva illac. Un \VUgras{soldato} \textsuperscript{(10)},
qui partiva in permission, habeva le complacentia de ceder me su vice, assi que
io habeva bastante tempore pro demandar alcun informationes al \VUgras{chef del
station} \textsuperscript{(13)}, qui conversava con le \VUgras{commissario}
\textsuperscript{(12)} del surveliantia e con le \VUgras{gendarme}
\textsuperscript{(11)} de servicio.

%%K t12-tit-1 sub
\VUpk{10.2pt}{40}{Inscription del bagages.}
\VUsaut{3.0mm}

%%K t12-05-1 p
5. --- Habente comprate un jornal al \VUgras{bibliotheca}
\textsuperscript{(14)}, io faceva pesar mi \VUgras{bagages}
\textsuperscript{(17)} que un \VUgras{portator} \textsuperscript{(16)} justo
habeva apportate per un \VUgras{carro de pacchettos} \textsuperscript{(19)};
\VUgras{le empleato} \textsuperscript{(22)} incargate del \VUgras{apparato de
pesar} \textsuperscript{(21)} me faceva saper, que mi cofro pesa trenta-cinque
kilogrammas; isto faceva un \VUgras{excesso de peso} de cinque kilogrammas; pro
le qual io debe un pagamento supplementari, al \VUgras{guichet} del bagages
\textsuperscript{(23)}.

%%K t12-06-1 p
6. --- Proque io esseva troppo tosto, io habeva tempore libere pro examinar le
circulation del viageros in le station. Unes deponeva o reprendeva lor bagages
de pauc valor al \VUgras{deposito (de bagages)} \textsuperscript{(25)}; alteres
consultava le \VUgras{horarios} \textsuperscript{(26)}. Le viageros arrivante
hastava al \VUgras{exito} \textsuperscript{(32)} con lor \VUgras{valise}
\textsuperscript{(24)} in le mano. Alcunes attendeva in le \VUgras{distribution
(de bagages)} \textsuperscript{(27)} e presentava lor \VUgras{billet}
\textsuperscript{(29)} pro reprender lor cofros, \VUgras{cassas}
\textsuperscript{(28)} o \VUgras{corbes} \textsuperscript{(30)}.

%%K t12-07-1 p
7. --- \VUgras{Le empleatos del accisa} \textsuperscript{(33)} curosemente
explorava le pacchettos pro constatar si on ha alco a declarar.

%%K t12-08-1 p
8. --- Alcun viageros iva al \VUgras{buffet} \textsuperscript{(34)}, ubi un
\VUgras{garson} \textsuperscript{(35)} se zelava a servir les. Illes debe hastar,
proque le \VUgras{traino} \textsuperscript{(36)}, que es un expresso, non
restara longemente in le station.

%%K t12-09-1 p
9. --- Postea io iva al trottoir de partita e cercava un \VUgras{vagon}
\textsuperscript{(37)} directe pro non esser obligate cambiar de vagon durante le
viage. Io ascendeva in un vagon a corridor que contine un \VUgras{compartimento}
\textsuperscript{(38)} pro fumatores e un compartimento reservate pro
« damas sol ».

%%K t12-tit-2 sub
\VUpk{10.2pt}{40}{Partita durante le nocte.}
\VUsaut{3.0mm}

%%K t12-10-1 p
10. --- Post haber ascendite le \VUgras{scabello} \textsuperscript{(40)}, cludite
le \VUgras{porta} \textsuperscript{(39)}, volvite le \VUgras{store}
\textsuperscript{(41)} e ponite mi parve bagages sur le \VUgras{rete}
\textsuperscript{(43)}, io sedeva sur le \VUgras{banchetta}
\textsuperscript{(42)}. Vidente le \VUgras{lampista} \textsuperscript{(44)}
accender le lampas, io pensava, que io debera passar un nocte in le vagon e io
appellava le empleato incargate del location del \VUgras{cossinos}
\textsuperscript{(45)} pro demandar un de illos.

%%K t12-11-1 p
11. --- Le conductor del traino veniva verificar le billets. Io le questionava
proque nos non ja habeva partite, proque il es le hora exacte. Ille me
respondeva que le \VUgras{chef del traino} \textsuperscript{(46)} se occupa de
facer poner bicyclettas in le \VUgras{furgon} \textsuperscript{(47)}. In plus,
on non habeva ancora cambiate tote le calefactores de pedes
\textsuperscript{(48)}.

%%K t12-12-1 p
12. --- Ma nos esseva tosto partiente, proque le \VUgras{fochista}
\textsuperscript{(50)}, sur su \VUgras{tender} \textsuperscript{(49)} prendeva
carbon pro excitar le foco del \VUgras{locomotiva} \textsuperscript{(54)}, e
\VUgras{le machinista} \textsuperscript{(51)} attendeva solmente le signal del
\VUgras{sub-chef del station} \textsuperscript{(52)}, qui esseva sur le
\VUgras{trottoir} \textsuperscript{(53)}.

%%K t12-13-1 p
13. --- Le \VUgras{caldiera} \textsuperscript{(56)} del locomotiva grunniva
sordemente; le \VUgras{camino} \textsuperscript{(55)} vomiva torrentes de fumo
miscite con \VUgras{vapor} \textsuperscript{(60)}. Un stridente \VUgras{sibilo}
\textsuperscript{(59)} sonava e le \VUgras{pistones} \textsuperscript{(58)}
comenciava mover se, impellente le \VUgras{rotas} \textsuperscript{(57)} del
pesante machina.

%%K t12-tit-3 sub
\VUsaut{3.0mm}
\VUpk{10.2pt}{40}{Partita!}
\VUsaut{3.0mm}

%%K t12-14-1 p
14. --- Le rapiditate del traino, currente sur le \VUgras{railes}
\textsuperscript{(64)} esseva accelerate; le \VUgras{trottoirs}
\textsuperscript{(65)} spareva; nos preterpassava le \VUgras{latrinas}
\textsuperscript{(66)}, e anque un serie de vagones garate sur le altere
\VUgras{via} \textsuperscript{(63)}: \VUgras{vagones-lectos}
\textsuperscript{(67)}, un \VUgras{vagon-restaurante} \textsuperscript{(68)}, un
\VUgras{vagon postal} \textsuperscript{(69)}.

%%K t12-noto-1 noto
\VUnotes{143.2mm}{%
\textit{(*) Nota.} --- \textit{Naturalmente le description del viage es secundo
le frontiera de ante le guerra.}%
}

%%K t12-15-1 p
15. --- Postea, le \VUgras{station del mercantias} \textsuperscript{(71)} passava
ante nostre oculos. Sub le hangares, empleatos cargava le \VUgras{mercantias}
\textsuperscript{(72)} expedite per traino lente. \VUgras{Homines de esquadra}
\textsuperscript{(76)} presso le \VUgras{reservoir de aqua}
\textsuperscript{(73)} manovrava \VUgras{vagones pro bestias}
\textsuperscript{(75)} e \VUgras{vagones a platteforma} \textsuperscript{(79)}
pro formar un \VUgras{traino de mercantias} \textsuperscript{(80)}.

%%K t12-16-1 p
16. --- Nos transversava le ultime \VUgras{bifurcation} \textsuperscript{(78)},
presso le qual stava le agulliero; nos preteriva le \VUgras{disco de signal}
\textsuperscript{(86)} e le station spareva in le lontano.

%%K t12-17-1 p
17. --- Tosto nos transversava un \VUgras{ponte de ferro} \textsuperscript{(74)}
que transpassa le \VUgras{fluvio} \textsuperscript{(81)}. Post transversar iste
ponte, nos sequeva le riviera, sur le qual potente \VUgras{remorcatores}
\textsuperscript{(82)} tirava pesante \VUgras{chalanas} \textsuperscript{(83)}.
Nos in plus videva un \VUgras{nave de viageros} \textsuperscript{(84)}, quando
illo abordava un \VUgras{pontone} \textsuperscript{(85)}.

%%K t12-18-1 p
18. --- Post preterpassar rapidissimemente plure stationes, nos transversava
alcun \VUgras{tunnels} \textsuperscript{(87)} e lassava detra nos
\VUgras{parve casas} innumerabile \textsuperscript{(88)} de \VUgras{guardas de
barriera}. Cunate per le succussas del traino, io profundemente addormiva.

%%K t12-tit-4 sub
\VUsaut{3.0mm}
\VUpk{10.2pt}{40}{Station de Deutsch-Avricourt.}
\VUsaut{3.0mm}

%%K t12-19-1 p
19. --- Io esseva subitemente evelliate per le voce de un empleato critante:
« Deutsch-Avricourt! \textsuperscript{(*)}. Tote le mundo descende! » Habeva
loco le inspection del dogana e tote le tediose formalitates del frontiera. Io
debeva conformar mi horologio de tasca al \VUgras{grande horologio}
\textsuperscript{(89)} del station, que esseva troppo avantiate per cinquanta-cinque
minutas; a pena evelliate, io distingueva difficilemente le \VUgras{grande
monstrator} \textsuperscript{(90)} del \VUgras{parve} \textsuperscript{(91)}; le
\VUgras{quadrante} \textsuperscript{(92)} mesme esseva toto confuse a causa del
nebuletta matinal.

%%K t12-20-1 p
20. --- Durante que le doganeros faceva lor inspection, io oscitante decifrava
le \VUgras{affiches} \textsuperscript{(93)} qui decora le muros del station. Io
dejunava pro passar le tempore, consultava le carta del \VUgras{rete}
\textsuperscript{(94)} german pro vider quante distantia ancora me separa de
Strasbourg, e io reentrava mi vagon, refortificate per le sperantia de tosto
attinger le scopo de mi viage.
\VUsaut{6.0mm}
\VUfilet{22mm}
